% Generated draft: requires mathematical and source review.
\documentclass[11pt]{article}
\usepackage[margin=25mm]{geometry}
\usepackage{fontspec}
\setmainfont{Latin Modern Roman}
\usepackage{amsmath,amssymb,mathtools,unicode-math}
\setmathfont{Latin Modern Math}
\usepackage{enumitem,needspace,xurl,hyperref}
\setlength{\parindent}{0pt}
\setlength{\parskip}{5pt}
\setlength{\emergencystretch}{3em}
\newcommand{\N}{\mathbb{N}}
\newcommand{\Z}{\mathbb{Z}}
\newcommand{\Q}{\mathbb{Q}}
\newcommand{\R}{\mathbb{R}}
\newcommand{\C}{\mathbb{C}}
\newcommand{\F}{\mathbb{F}}
\newcommand{\A}{\mathbb{A}}
\newcommand{\PP}{\mathbb{P}}
\newcommand{\E}{\mathbb{E}}
\newcommand{\eps}{\varepsilon}
\newcommand{\dd}{\,\mathrm{d}}
\begin{document}
\section*{UnsolvedMath ID: 800277}
\textbf{Source problem code:} catalogue entry
\subsection{Definitions and mathematical statement}
Interpret the source's projective plane as the complex projective plane \(\mathbb{P}^{2}_{\C}\). Let \(C\subset\mathbb{P}^{2}_{\C}\) be a smooth plane cubic: the zero locus of a homogeneous polynomial of degree \(3\), with no singular points. Put
\[
X_C=\mathbb{P}^{2}_{\C}\setminus C.
\]
A holomorphic immersion into \(\C^2\) means a holomorphic map \(f:X_C\to\C^2\) whose differential
\[
\dd f_x:T_xX_C\longrightarrow T_{f(x)}\C^2
\]
is injective at every \(x\in X_C\). Since both tangent spaces have complex dimension \(2\), this differential must be an isomorphism. Global injectivity of \(f\) is not required.

The question is whether, for every smooth plane cubic \(C\subset\mathbb{P}^{2}_{\C}\), there exists such a map:
\[
\forall C\text{ smooth plane cubic},\quad
\exists f:X_C\to\C^2\text{ holomorphic},\quad
\forall x\in X_C,\quad \operatorname{rank}_{\C}(\dd f_x)=2.
\]

As background, the source asserts that \(X_C\) is algebraically parallelizable, meaning that its algebraic tangent bundle is isomorphic to \(X_C\times\C^2\), and that it cannot be algebraically immersed in \(\C^2\). These are supplied background assertions, not conclusions established here; the source does not elaborate its terminology for algebraic immersion. The source also states that the holomorphic immersion exists for some smooth cubics, without identifying those cubics or providing a criterion. Thus the requested existence question concerns an arbitrary smooth cubic, beyond the supplied affirmative special cases.
\subsection{Short English statement}
Does the complement of every smooth complex plane cubic admit a holomorphic map to \(\C^2\) with everywhere injective differential? The source reports affirmative examples for some smooth cubics and contrasts this with an asserted obstruction to algebraic immersion. Its supplied status is open; this draft does not independently verify that status.
\subsection{Sources}
\begin{itemize}
\item Local UnsolvedMath registry: permanent ID \(800277\), exact source problem code \texttt{catalogue entry}, title \emph{Several Complex Variables Problem List --- Problem 77}. The supplied record labels the problem open and was updated on 2026-10-16.
\item Buzzard, Fornaess, and Krantz, \emph{The Several Complex Variables Problem List} (1995), Problem 77: \url{https://arxiv.org/abs/math/9509202}. The registry records access on 2026-10-16 and extraction from arXiv LaTeX; these are supplied provenance details, not a claim of independent access or verification.
\item The problem text attributes the question to R. Narasimhan, Department of Mathematics, University of Chicago, Chicago, IL 60637. The registry separately lists Buzzard, Fornaess, and Krantz under proposer attribution, with proposed year 1995.
\item The source directs readers to its section entitled ``Stein manifolds'' for related problems. No specific related problem numbers or references are supplied.
\end{itemize}
\end{document}
